\documentclass[conference,compsoc]{IEEEtran}

\usepackage{booktabs}
\usepackage{url}
\usepackage[hidelinks]{hyperref}

% Generated by scripts/generate_paper_tables.py; do not hand edit.
\newcommand{\benchmarkname}{AuthZBench-SaaS}
\newcommand{\releaseversion}{v1-prep draft}
\newcommand{\publictaskcount}{54}
\newcommand{\publicappcount}{6}
\newcommand{\vulnerabletaskcount}{21}
\newcommand{\securecontrolcount}{33}
\newcommand{\denialcontrolcount}{19}
\newcommand{\authorizedallowcount}{14}


\begin{document}

\title{\benchmarkname: Evaluating Replayable Authorization-Failure Proofs in AI Security Agents}

\author{\IEEEauthorblockN{Anonymous Authors}}

\maketitle

\begin{abstract}
AI security agents can produce plausible vulnerability reports without proving
that an authorization failure occurred. \benchmarkname{} \releaseversion{}
introduces a focused benchmark for this proof gap. The benchmark contains
\publicappcount{} synthetic SaaS targets and \publictaskcount{} public tasks
covering tenant, organization, object, role, API-token, scope, sharing, and
admin-action boundaries. It evaluates whether an agent can submit
backend-replayable evidence for vulnerable tasks while avoiding false reports
on secure controls. The current public split includes \vulnerabletaskcount{}
vulnerable tasks, \securecontrolcount{} secure-control tasks,
\denialcontrolcount{} denial controls, and \authorizedallowcount{}
authorized-allow controls. This paper frames the released \releaseversion{}
repository as a research artifact and methodology foundation, not as a hosted
leaderboard or v1 community benchmark.
\end{abstract}

\section{Introduction}

SaaS authorization bugs are difficult to evaluate with prose alone. A useful
agent must identify the actor, tenant or organization, object, role, token, or
scope boundary and provide backend evidence that can be replayed by the scorer.
\benchmarkname{} focuses on this narrow capability: proving authorization
failures while staying quiet on secure controls.

\section{Benchmark Design}

\benchmarkname{} uses synthetic local SaaS fixtures so scorer-owned replay can
be deterministic and public users can reproduce the public split. The benchmark
does not claim production vulnerability-discovery capability, broad cyber
capability, or hosted leaderboard readiness.

\begin{table}[t]
\centering
\caption{Public task mix by target family.}
% Generated by scripts/generate_paper_tables.py; do not hand edit.
\scriptsize
\begin{tabular}{@{}lrrrrrr@{}}
\toprule
Family & Tasks & Vuln. & Ctrl. & Deny & Allow & Flow \\
\midrule
api tokens & 10 & 4 & 6 & 3 & 3 & 1 \\
audit settings & 9 & 4 & 5 & 2 & 3 & 1 \\
billing & 13 & 5 & 8 & 5 & 3 & 1 \\
file sharing & 8 & 3 & 5 & 3 & 2 & 0 \\
project mgmt & 9 & 4 & 5 & 3 & 2 & 1 \\
support & 11 & 4 & 7 & 5 & 2 & 1 \\
\midrule
\textbf{Total} & 60 & 24 & 36 & 21 & 15 & 5 \\
\bottomrule
\end{tabular}

\end{table}

\section{Scoring}

Vulnerable tasks require replayable exploit proof, correct boundary reasoning,
control replay, and safe behavior for a full pass. Secure controls require no
findings. Release-facing interpretation separates exploit proof, boundary
reasoning, false positives, control execution, authorized-allow behavior, and
live-target request coverage.

\section{Experimental Results}

Current public baselines are diagnostic public-split evidence. They show that
agents can sometimes produce replayable exploit evidence while still failing
full vulnerable-task passes because boundary reasoning remains weak.

\begin{table*}[t]
\centering
\caption{Current public baseline evidence generated from repository artifacts.}
% Generated by scripts/generate_paper_tables.py; do not hand edit.
\footnotesize
\begin{tabular}{@{}llrrrrrrr@{}}
\toprule
Baseline & Harness & Tasks & V0 pass & Exploit & Boundary & FP & Allow & Requests \\
\midrule
Scripted sanity baseline (v0.0) & scripted & 46 & 46 & 1 & 1 & 0 & 1 & n/a \\
qwen3-coder-next run 1 (v0.0) & no-tools-model & 46 & 27 & 0 & 0 & 0 & 1 & n/a \\
qwen3-coder-next run 2 (v0.0) & no-tools-model & 46 & 27 & 0.0526 & 0 & 0 & 1 & n/a \\
claude-haiku-4.5 run 1 (v0.0) & no-tools-model & 46 & 26 & 0.2632 & 0 & 0.037 & 1 & n/a \\
claude-haiku-4.5 run 2 (v0.0) & no-tools-model & 46 & 27 & 0.0526 & 0 & 0 & 1 & n/a \\
claude-sonnet-4.6 run 1 (v0.0) & no-tools-model & 46 & 27 & 0.6316 & 0 & 0 & 1 & n/a \\
claude-sonnet-4.6 run 2 (v0.0) & no-tools-model & 46 & 26 & 0.4211 & 0 & 0.037 & 1 & n/a \\
glm-5 run 1 (v0.0) & no-tools-model & 46 & 27 & 0.2105 & 0 & 0 & 1 & n/a \\
glm-5 run 2 (v0.0) & no-tools-model & 46 & 27 & 0.0526 & 0 & 0 & 1 & n/a \\
claude-sonnet-4.6 run 1 (v0.0) & tool-agent & 46 & 27 & 0.7368 & 0 & 0 & 1 & 1 \\
claude-sonnet-4.6 run 2 (v0.0) & tool-agent & 46 & 27 & 0.7368 & 0 & 0 & 1 & 1 \\
\bottomrule
\end{tabular}

\end{table*}

\section{Related Work}

\benchmarkname{} sits at the intersection of API authorization testing,
cybersecurity evaluation for language models, agent benchmarks, and benchmark
documentation. OWASP API Security Top 10 distinguishes object-level and
function-level authorization risks~\cite{owasp-api-top10-2023}. Cybersecurity
evaluation suites such as CyberSecEval motivate narrower, evidence-backed cyber
capabilities~\cite{bhatt2024cyberseceval2,bhatt2024cyberseceval3}. Agent
benchmarks such as SWE-bench motivate task-level environments and reproducible
grading~\cite{jimenez2024swebench}. Model cards and datasheets inform the
repository's claim ledger and documentation
discipline~\cite{mitchell2019modelcards,gebru2021datasheets}. The IEEE S\&P
scaffold follows the venue-facing formatting path noted in the current call for
papers~\cite{ieeesp2026cfp}.

\section{Ethical Considerations}

The targets are synthetic local fixtures and should not be exposed to the
public internet. Public tasks are inspectable and support reproducibility; they
must not be treated as private leaderboard tasks. Private holdout bodies,
routes, seeds, and raw evidence remain out of public Git history.

\section{Open Science and Artifact Availability}

The public repository contains task manifests, scorer code, validation scripts,
public-safe charts, baseline summaries, and release evidence. Maintainer-only
private holdout evidence is exposed only through aggregate redacted summaries.
The artifact packet under \texttt{artifact/} documents the intended public
validation path.

\section{Limitations}

The \releaseversion{} public split is small, synthetic, and inspectable. The
current baseline table is not a definitive model ranking. Hosted leaderboard
operation, rotating private holdouts, broader isolation, external review, and
100+ task scale are future v1 work.

\section{LLM Usage Considerations}

Because the benchmark evaluates model and agent behavior, any paper or report
using \benchmarkname{} should disclose model labels, harness type, tool access,
task split, benchmark fingerprint, repeated-run provenance, and known
automation assistance used to prepare benchmark artifacts.

\section{Evidence Map}

\begin{table}[t]
\centering
\caption{Claim evidence map generated from repository documentation.}
% Generated by scripts/generate_paper_tables.py; do not hand edit.
\scriptsize
\begin{tabular}{@{}p{0.26\columnwidth}p{0.17\columnwidth}p{0.47\columnwidth}@{}}
\toprule
Area & Status & Generated evidence \\
\midrule
Public split & Supported & 54 tasks; 21 vulnerable; 33 controls \\
V0.0 release snapshot baselines & Supported & 6 registry entries; current v1 reruns pending \\
Private leaderboard-candidate rows & Partially supported & 1 eligible of 2 tracked submissions; not hosted leaderboard operation \\
Evidence-map claim ledger & Mixed & 14 supported; 6 partial; 8 unsupported rows \\
\bottomrule
\end{tabular}

\end{table}

\section{Conclusion}

\benchmarkname{} \releaseversion{} is best understood as an evidence-backed
research artifact for studying SaaS authorization proof quality. The next step
is not stronger release wording, but independent review, variance analysis,
artifact packaging, and v1 task expansion.

\bibliographystyle{IEEEtran}
\bibliography{references}

\end{document}
